\documentclass[11pt]{article}

\usepackage[a4paper,margin=1in]{geometry}
\usepackage{amsmath,amssymb,amsthm,mathtools}
\usepackage{microtype}

\newtheorem{theorem}{Theorem}[section]

\newcommand{\qbinom}[2]{\genfrac{[}{]}{0pt}{}{#1}{#2}}

\title{A Pairing Index--Width Theorem for the Partition Function}
\author{Chung-Ming Liung, Yen-Tarng Wang, and Chih-Yi Tsay}
\date{June 25, 2026}

\begin{document}
\begin{theorem}[Pairing index--width theorem]\label{thm:pairing-index-width}
For every $n\ge0$ and $a,b\ge0$, the number of partitions of $n$ having pairing index $a$ and pairing width $b$ equals the number of partitions of $n$ having exactly $a$ parts and largest part $b$.
\end{theorem}

\begin{proof}
For $B\ge0$, we put
\[
        F_B(x,q)=\sum_{W(\lambda)\le B}x^{T(\lambda)}q^{|\lambda|}.
\]
We prove
\begin{equation}\label{eq:pairing-index-width}
        F_B(x,q)=\frac{1}{(xq;q)_B}.
\end{equation}
The case $B=0$ is immediate.  We assume $B\ge1$.

After maximal pairing, we write $\lambda$ uniquely as a paired portion $\pi$ and a distinct unpaired portion $\delta$: for each extracted pair $j+j$, we place one copy of $j$ in $\pi$, while $\delta$ consists of the parts left unpaired.  Then
\[
        |\lambda|=2|\pi|+|\delta|,
        \qquad
        T(\lambda)=\pi_1+s(\delta),
\]
with $\pi_1=0$ when $\pi$ is empty.  Moreover, the pair span is $2\ell(\pi)$.

We set
\[
        m=\left\lfloor\frac{B}{2}\right\rfloor,
        \qquad
        h=\left\lfloor\frac{B-1}{2}\right\rfloor.
\]
The condition $W(\lambda)\le B$ is equivalent to the two independent conditions
\[
        \ell(\pi)\le m,
        \qquad
        \delta=\varnothing \text{ or } \delta_1-s(\delta)\le h.
\]
Consequently $F_B$ factors into a paired and an unpaired contribution.  For the paired contribution, conjugating $\pi$ gives
\[
        \sum_{\ell(\pi)\le m}x^{\pi_1}q^{2|\pi|}
        =\frac{1}{(xq^2;q^2)_m}.
\]
By Lemma~, the unpaired contribution is
\[
        \frac{1}{(xq;q^2)_{h+1}}.
\]
Hence
\[
        F_B(x,q)
        =\frac{1}{(xq^2;q^2)_{\lfloor B/2\rfloor}(xq;q^2)_{\lfloor B/2\rfloor}}
        =\frac{1}{(xq;q)_B},
\]
which proves .

The right-hand side is the generating function for partitions whose largest part is at most $B$, with $x$ marking the number of parts.  Thus
\[
        \#\{\lambda\vdash n:T(\lambda)=a,\ W(\lambda)\le B\}
        =\#\{\mu\vdash n:\ell(\mu)=a,\ \mu_1\le B\}.
\]
Subtracting the corresponding identity with $B-1$ in place of $B$ gives the assertion with pairing width and largest part both equal to $B$.
\end{proof}
\end{document}
